\documentclass[11pt]{article}

% Compiles with pdflatex (Overleaf: default). Only ubiquitous packages.
\usepackage[margin=1in]{geometry}
\usepackage{amsmath,amssymb}
\usepackage{booktabs}
\usepackage{enumitem}
\usepackage{xcolor}
\usepackage[hidelinks]{hyperref}
\usepackage{microtype}

\newcommand{\code}[1]{\texttt{#1}}
\newcommand{\layer}[1]{\textsc{#1}}

\title{\textbf{verifiable-llm}: Practical, Precisely-Scoped Provenance\\
for Local Language-Model Inference}

\author{Oscar Tiz\\ % TODO: set your name for submission
\texttt{github.com/oscartiz/verifiable-llm}}

\date{July 2026}

\begin{document}
\maketitle

\begin{abstract}
When a party claims ``model $M$ produced text $T$ from prompt $P$,'' a verifier
normally has no way to check it: they cannot tell which weights ran, or whether
the output was edited after the fact. Proving a full transformer forward pass in
zero knowledge would settle the question but is infeasible on commodity
hardware. We present \textbf{verifiable-llm}, a system that instead composes
three layers of increasing cryptographic strength, each with an explicit and
deliberately understated security statement. \layer{Layer~1} binds a claim to a
specific $(\text{model}, \text{prompt}, \text{sampler}, \text{per-step logits},
\text{tokens})$ tuple using a \textsc{blake3} Merkle commitment over the
quantized weights and a domain-separated hash chain over decode steps.
\layer{Layer~2} adds probabilistic soundness: Fiat--Shamir spot-check challenges
that re-execute random blocks of the committed computation on CPU, catching a
cheating prover with tunable probability. \layer{Layer~3} proves, in zero
knowledge, that the emitted greedy token is the argmax of the committed logit
vector, via a transparent STARK (fast; \emph{not} formally zero-knowledge) and a
halo2 inner-product-argument variant (formally zero-knowledge; $O(\text{vocab})$
to prove). A distinguishing feature is that the system quantifies an attack
against \emph{its own} float-tolerance verifier---sub-tolerance perturbations are
amplified $30$--$80\times$ by the network and can steer decoding---and closes it
with a deterministic, exactly-verifiable backend. We report measured costs on an
Apple~M4 laptop: a decode-argmax STARK proves in ${\approx}0.7$\,s and verifies
in ${<}1$\,ms with ${\approx}78$\,KB proofs. The system is implemented in Rust
($6$ crates) with adversarial tests that catch a cheating prover at every layer.
\end{abstract}

\section{Introduction}

Text, code, and decisions increasingly come from large language models (LLMs)
run locally on consumer hardware. This raises a provenance question with legal,
scientific, and safety stakes: given an output, can a third party later verify
\emph{which} model produced it and that the output was not fabricated or edited?
Today the answer is essentially ``no''---the claimant controls the machine, and
nothing binds the output to a computation.

The cryptographic ideal is \emph{zero-knowledge machine learning} (zkML): prove
in zero knowledge that $y = M(x)$ for a committed model $M$. For a
billion-parameter transformer this is currently impractical on a laptop: the
arithmetic circuit is enormous, and hiding it entirely is the wrong cost/benefit
trade for most provenance needs. Rather than pay for the maximal guarantee,
\textbf{verifiable-llm} decomposes the problem into layers whose guarantees a
practitioner can compose and reason about, and states each guarantee with the
precision an applied cryptographer would demand---including what it does
\emph{not} provide.

\paragraph{Contributions.}
\begin{itemize}[leftmargin=1.2em,itemsep=2pt]
\item A layered provenance protocol for local quantized-llama inference:
binding commitments (\S\ref{sec:l1}), interactive spot checks with a quantified
catch probability (\S\ref{sec:l2}), and zero-knowledge decode proofs
(\S\ref{sec:l3}), each with a stated adversary and non-goals.
\item An adversarial analysis of tolerance-based verification
(\S\ref{sec:drift}): we show a prover running the \emph{real} model but
perturbing every committed activation below the verifier's tolerance can steer
greedy decoding, because the network amplifies sub-tolerance drift
$30$--$80\times$. We close this with a deterministic, bit-exact backend whose
challenges verify with \emph{zero} tolerance.
\item Two interchangeable decode-argmax proof backends behind one commitment
format (\S\ref{sec:l3}): a fast transparent STARK that is a succinct argument
but not formally zero-knowledge, and a halo2 IPA circuit that is formally
zero-knowledge at $O(\text{vocab})$ cost. We make the honest limitation of each
explicit.
\item An open-source Rust implementation and an evaluation on commodity hardware
(\S\ref{sec:eval}), with reproducible proof-cost benchmarks and an adversarial
test suite that catches a cheating prover at each layer.
\end{itemize}

\section{Threat Model and Design Goals}\label{sec:threat}

\paragraph{Adversary.} The adversary is a \emph{prover}: a party who ran (or
claims to have run) some computation and later wants to lie about it. Concretely
the prover may (i) substitute, fine-tune, or re-quantize the model; (ii)
fabricate logits or tokens that the model never produced; (iii) edit the
transcript after generation; or (iv) attempt to answer verification challenges
inconsistently with the committed computation. The verifier is honest but may be
a \emph{different machine}, possibly without a GPU, that only exchanges files
with the prover.

\paragraph{Explicit non-goals.} We do \emph{not} prove a full forward pass in
zero knowledge; we do not defend against a prover who honestly runs the correct
model (that is the point of provenance, not an attack); and we do not certify
model quality or safety. Sampler support for the strongest layer is greedy
(argmax) decoding; top-$p$/temperature proofs are future work.

\paragraph{Cryptographic assumptions.} Collision resistance of
\textsc{blake3}~\cite{blake3}; the Fiat--Shamir heuristic in the random-oracle
model~\cite{fiatshamir}; the conjectured soundness of the FRI-based
STARK~\cite{stark} ($\approx\!100$-bit) and of the halo2 inner-product
argument~\cite{bulletproofs,halo} over the Pasta curves. The deterministic
backend additionally relies only on IEEE-754 being fully specified for the basic
operations it uses.

\section{Layer 1: Binding Commitments}\label{sec:l1}

Layer~1 makes a claim \emph{binding}, not \emph{correct}: it pins exactly what
was claimed so it cannot be rewritten, without yet asserting the computation was
performed honestly.

\paragraph{Model commitment.} We build a Merkle tree over the tensors of the
GGUF weight file. Leaf $i$ is
$\text{H}_{\text{leaf}}(\text{name}\,\|\,\text{ggml\_type}\,\|\,\text{shape}\,\|\,\text{raw quantized bytes})$,
computed over the bytes \emph{exactly as stored on disk}---no dequantization, so
nothing float-ambiguous enters the model commitment. Leaves are sorted by tensor
name (so the root is invariant to on-disk order, which re-quantization tools may
change) and combined with the RFC~6962 tree shape~\cite{rfc6962} using
domain-separated leaf and node hashing (\code{vllm/merkle-leaf/v1},
\code{vllm/merkle-node/v1}). Domain separation blocks second-preimage and
leaf/internal-node confusion attacks; a single-leaf tree hashes under the leaf
domain only, so it can never be reinterpreted as an interior node. Two files
with the same root are byte-identical in every committed payload, name, type,
and shape up to \textsc{blake3} collision resistance.

\paragraph{Transcript hash chain.} Generation is bound by a running chain:
\begin{align*}
\text{chain}_0 &= \text{H}(\text{model\_root} \,\|\, \text{prompt\_tokens} \,\|\, \text{sampler\_cfg}),\\
\text{chain}_i &= \text{H}(\text{chain}_{i-1} \,\|\, i \,\|\, \text{token}_i \,\|\, \text{H}(\text{quantize}(\text{logits}_i))).
\end{align*}
Every variable-length field is length-prefixed and every hash input is
domain-tagged; the sampler configuration is committed as raw IEEE-754 bits (mode,
temperature, top-$p$, RNG seed), avoiding JSON-canonicalization ambiguity.
Logits are committed at fixed point, $q = \mathrm{round}(x \cdot 2^{16})$ as
\code{i32}: coarse enough to hash stably on one backend, fine enough to leave the
argmax and any downstream circuit unaffected (\S\ref{sec:l3}). \code{NaN} logits
abort rather than being silently committed; $\pm\infty$ saturates. The final
chain value binds the prover to one ordered tuple; any retroactive edit changes
it, detectable by re-running the chain with no model or GPU.

\paragraph{What Layer 1 does not do.} It does not prove the logits came from
running $M$: a prover could commit to fabricated logits. Establishing that the
committed computation actually follows from the committed weights is the job of
Layer~2.

\section{Layer 2: Interactive Spot Checks}\label{sec:l2}

Generating with tracing additionally commits every intermediate activation
\emph{cell} $(p, j)$---the hidden state entering block $j$ at position $p$, plus
the state exiting the last block---quantized and folded into a Merkle tree whose
root is bound into the chain as generation runs. A verifier then challenges a
random subset and re-executes it.

\paragraph{Challenge derivation.} The verifier (or a Fiat--Shamir transform)
draws $k$ distinct cells from
\[
\text{XOF}\big(\text{final\_chain}\,\|\,\text{trace\_root}\,\|\,k\,\|\,\text{nonce}\,\|\,\text{space params}\big),
\]
where the space parameters (positions, layers, first-logit position) are
absorbed so the challenge space cannot be reshaped. Sampling is without
replacement.

\paragraph{Response and check.} For each challenged block the prover reveals the
block's committed input cells at positions $0..p$ and its output at $p$, each
with a Merkle inclusion proof; the verifier re-executes the block with the
committed weights over the revealed inputs and compares. Head challenges
additionally reveal the step's quantized logits, which must hash \emph{exactly}
to the chain's committed logits digest, and layer-$0$ inputs must equal the
committed tokens' embeddings. The entire verifier runs on CPU, so a second
machine without a GPU can check the claim.

\paragraph{Soundness.} Call a \emph{cheat} a committed cell inconsistent with the
committed model's computation on the committed inputs. If a fraction $f$ of the
$N$ committed cells are cheats, then $k$ distinct challenges miss all of them
with probability at most $(1-f)^k$; equivalently
\[
k = \left\lceil \frac{\ln \delta}{\ln(1-f)} \right\rceil
\]
challenges catch the prover with probability at least $1-\delta$
(Table~\ref{tab:catch}). Sampling without replacement makes the bound tighter
than the with-replacement $(1-f)^k$ it is stated against.

\begin{table}[t]
\centering
\small
\begin{tabular}{lccc}
\toprule
cheat fraction $f$ & $k$ (95\%) & $k$ (99\%) & $k$ (99.9\%)\\
\midrule
20\% & 14 & 21 & 31\\
10\% & 29 & 44 & 66\\
5\%  & 59 & 90 & 135\\
1\%  & 299 & 459 & 688\\
\bottomrule
\end{tabular}
\caption{Distinct challenges needed to catch a prover cheating on a fraction $f$
of committed cells, at the exact $\lceil \ln\delta / \ln(1-f)\rceil$ bound.}
\label{tab:catch}
\end{table}

\paragraph{Grinding.} With a pure Fiat--Shamir (empty) nonce a prover can
re-generate until the derived challenges miss its cheats, at expected cost
$(1-f)^{-k}$ generations. A verifier-chosen nonce, issued \emph{after} the
transcript is committed, removes this entirely; the prover's \code{respond}
refuses any challenge that is not the exact derivation for the committed
$(\text{transcript}, \text{nonce})$.

\section{The Bounded-Drift Attack and Deterministic Verification}\label{sec:drift}

Layer~2's re-execution cannot use exact equality on the default path.
Cross-backend logit reproduction is impossible at any usable precision: on
Llama-3.2-1B~\cite{llama3} Q4\_K\_M we measure a maximum absolute CPU-versus-Metal
logit deviation of ${\approx}0.49$ (mean ${\approx}0.08$) over the $128{,}256$-dim
vector, from differing quantized matmul kernels; the argmax agrees but hash
equality fails at every precision. The verifier therefore accepts a re-executed
block within a per-element tolerance $\tau$.

This opens an escape hatch, which we quantify against \emph{our own} system. An
adversary who runs the real model but adds a perturbation $\delta$ with
$\|\delta\|_\infty \le \tau$ at every committed cell produces a trace that is
self-consistent up to $\tau$ and passes every per-element challenge by
construction. The question is what that buys at the output. We find:

\begin{itemize}[leftmargin=1.2em,itemsep=2pt]
\item \textbf{Amplification.} A single $\pm\tau$ injection at one hook displaces
the final logits by $30$--$80\times\,\tau$ (Jacobian amplification through the
remaining blocks), saturating around $40$--$70\times$ at $\tau = 0.5$.
\item \textbf{Token steering.} Against an honest top-1/top-2 logit gap with
$10$th percentile ${\approx}0.33$, random full-stack $\pm\tau$ drift flips the
first token within ${\approx}10$ steps even at $\tau = 0.01$---below the honest
cross-backend noise floor itself. A targeted last-layer perturbation
$\delta = \tau\cdot\mathrm{sign}(w_b - w_a)$ (from the top-2 LM-head rows) flips
the argmax of $24$ of $48$ steps at $\tau = 0.01$.
\end{itemize}

Per-element tolerance therefore cannot guarantee token provenance for a float
transformer: this is a structural limit, not a parameter-tuning bug. A shipped
distributional check (a per-cell \emph{mean}-deviation bound, exploiting that
honest drift concentrates two orders of magnitude below its own max) narrows the
attack ${\approx}10\times$ but does not close it, and we say so.

\paragraph{The fix is determinism.} We add a fixed-evaluation-order CPU backend
(\code{det-cpu-v1}) that is bit-identical across runs and machines. It requires
exactly three things: (i) fixed reduction order (sequential scalar loops;
parallelism only over independent outputs), so IEEE-754's exact specification for
$+,-,\times,\div,\sqrt{}$ makes the result deterministic; (ii) no \code{libm}
(fixed-order polynomial \code{exp}/\code{sin}/\code{cos}, ${\approx}10^{-13}$
relative error); and (iii) requantization of the hidden state to the trace grid
at every hook, so a committed cell \emph{is} the computation state and the
verifier's recomputation is bit-identical. Verification then uses exact
\code{i32} equality with \emph{zero} tolerance: a single-quantum ($2^{-8}$)
perturbation on one element of one cell---undetectable at any float
tolerance---is caught. This closes the bounded-drift attack for deterministic
transcripts, at ${\approx}8\times$ lower generation throughput and
${\approx}5$\,GB extra memory (weights dequantized to \code{f32} up front).

\section{Layer 3: Zero-Knowledge Decode Proofs}\label{sec:l3}

Layers~1--2 establish that the committed logits are the model's true logits
(exactly, on the deterministic path; within an accumulated drift bound on the
float path). Layer~3 proves the emitted token was chosen correctly from those
logits, without revealing them. Both backends prove the same statement about one
greedy step:
\begin{quote}\itshape
For a public token index $c$ and a chain-bound digest $d$, the prover knows a
logit vector $x$ and salt $s$ such that $\mathrm{Commit}_s(x) = d$ and
$x[c] \ge x[i]$ for every $i$.
\end{quote}

\paragraph{Argmax by range check.} Both circuits introduce a private
claimed-maximum $m$ and enforce $m - x[i] \in [0, 2^{27})$ for all $i$ via a
$27$-bit decomposition, together with $x[c] = m$ selected at the public index.
Because $2^{27}$ is far below the field modulus, the field difference equals the
true signed-integer difference with no wraparound; a non-maximal claim forces
some $m - x[j] < 0$, i.e.\ a field element ${\approx}\!p$, which no $27$-bit
witness can represent, so the proof is infeasible. The bound caps the provable
logit spread at $2^{27}/2^{16} = 2048$ units, comfortably above observed spreads
(${\approx}40$).

\subsection{STARK backend (winterfell): the fast default}
The commitment is a salted Rescue-Prime~\cite{rescue} sponge over the Goldilocks
field, with the salt as a private capacity IV. An algebraic intermediate
representation (AIR) recomputes the sponge from the witness logits \emph{inside}
the proof (one Rescue round per row, one logit absorbed per row on average) and
simultaneously enforces the argmax constraints above via a boolean selector
column that must sum to one and may fire only at row $c$.

\emph{Security, stated precisely.} This is a succinct, \emph{transparent}
argument of knowledge (no trusted setup) with ${\approx}100$-bit conjectured
security, and the commitment is binding. It is \textbf{not} formal
zero-knowledge: the STARK library has no trace randomization, so each proof
reveals the trace polynomials' evaluations at the ${\approx}27$ FRI query points
of a coset domain---equivalently ${\approx}27$ random linear projections of the
$(\text{salt}\,\|\,\text{logits})$ vector. No individual logit is recoverable
(the system is underdetermined by ${\approx}5$ orders of magnitude, and the salt
blocks digest-level guessing), but a party already holding a candidate for the
\emph{entire} logit vector could confirm it. We call this reconstruction-safe but
confirmation-unsafe.

\subsection{halo2 backend: formal zero-knowledge}
The commitment is a salted Poseidon~\cite{poseidon} hash chain
$\text{acc}_{i+1} = \mathrm{Poseidon2}(\text{acc}_i, x_i)$, $\text{acc}_0 = s$,
using the same primitive the in-circuit gadget realizes (native/in-circuit
parity is pinned by tests). A running one-hot selector proves the argmax without
a per-token key: four running accumulators force exactly one selected row, at the
public index $c$ read from the instance column via a copy constraint, with the
selected value equal to a broadcast maximum---so a single verifying key serves
every token. The prover is halo2's inner-product argument over the Pasta curves:
\emph{transparent} (no trusted setup, like the STARK) \emph{and} it blinds every
committed witness polynomial, so the proof reveals nothing about the logits
beyond the public claim.

\emph{The cost.} The commitment is an in-circuit Poseidon over the whole vocab,
$O(\text{vocab})$ rows. This is affordable for small vocabularies but, projected
to a real llama vocabulary of $128{,}256$, requires a ${\approx}2^{23}$-row
circuit---minutes of proving and multiple GB. The STARK therefore remains the
fast default; halo2 is the opt-in path when \emph{formal} zero-knowledge is
required. Crucially, both backends fold a salted $32$-byte digest into the same
chain, so the transcript format is backend-agnostic and unchanged.

\paragraph{Composition.} Combining the layers: the emitted token is the exact
argmax of logits that are the committed model's true logits on the committed
inputs---exactly on the deterministic path, and within the accumulated drift
bound on the float path.

\section{Implementation}\label{sec:impl}

\textbf{verifiable-llm} is ${\approx}5{,}800$ lines of Rust across six crates.
\code{vllm-core} (GGUF parsing, Merkle tree, hash chain, transcript/trace/challenge
formats) depends only on \textsc{blake3} and \code{serde}, so a verifier builds
it anywhere without a GPU or ML framework. \code{vllm-infer} vendors candle's
quantized-llama with commitment hooks and the deterministic backend;
\code{vllm-verify} is the CPU re-execution verifier; \code{vllm-zk} and
\code{vllm-zk-halo2} are the two Layer-3 backends, isolated so their heavy
proof-system dependencies do not affect the default build; and \code{vllm-cli}
is the \code{vllm} binary. Attacker-controlled inputs (transcript, challenge,
response, and proof JSON) are parsed fail-closed: malformed input yields a clean
error, never a panic. Integration tests construct a tiny in-memory GGUF, so
continuous integration never downloads weights, and include a cheating prover
that is \emph{caught} at each layer: a self-consistent but incorrect trace
(caught by re-execution), sub-tolerance drift (caught by the mean check), a
single quantum under determinism (caught by exact equality), and forged Merkle
proofs, edited transcripts, wrong-salt digests, and mistagged decode proofs (all
rejected).

\section{Evaluation}\label{sec:eval}

We measure on an Apple~M4 MacBook~Air (16\,GB, macOS~26.5.1, \code{rustc}~1.95.0,
release build). Proof-cost benchmarks use synthetic logits and need no model
download; model-dependent numbers use Llama-3.2-1B-Instruct Q4\_K\_M.

\paragraph{Layer-3 proof costs.} Table~\ref{tab:l3} contrasts the two backends.
The STARK proves a $128{,}256$-vocab decode step in ${\approx}0.7$\,s and
verifies in under a millisecond with ${\approx}78$\,KB proofs; the native
commitment costs ${\approx}45$--$80$\,ms/token (machine-dependent) and is
opt-in. The halo2 backend (Table~\ref{tab:halo2}) is formally zero-knowledge but
$O(\text{vocab})$: at small vocabularies its \emph{cached} prove (reusing the
deterministic, cacheable parameters and keys) is $0.4$--$3.0$\,s with a flat
${\approx}5$\,KB proof, but its verification is linear in circuit size and it
does not scale to a full vocabulary in seconds.

\begin{table}[t]
\centering
\small
\begin{tabular}{lcc}
\toprule
 & STARK (\code{vllm-zk}) & halo2 (\code{vllm-zk-halo2})\\
\midrule
setup & none (transparent) & none (transparent)\\
formal ZK & no & \textbf{yes}\\
prove (vocab 128k) & ${\approx}0.7$\,s & minutes (est.)\\
verify & ${\approx}0.6$\,ms & seconds (est.)\\
proof size & $76$--$79$\,KiB & ${\approx}5$\,KB\\
commit / token & ${\approx}45$--$80$\,ms & $O(\text{vocab})$\\
\bottomrule
\end{tabular}
\caption{Decode-argmax proof backends. Both are transparent; they trade formal
zero-knowledge against proving cost at a real vocabulary.}
\label{tab:l3}
\end{table}

\begin{table}[t]
\centering
\small
\begin{tabular}{rccccc}
\toprule
vocab & $k$ & prove (cold) & prove (cached) & verify (cached) & proof\\
\midrule
32  & 12 & $1.30$\,s  & $0.42$\,s & $17.5$\,ms & 5\,KB\\
128 & 14 & $5.46$\,s  & $1.50$\,s & $19.6$\,ms & 5\,KB\\
256 & 15 & $10.96$\,s & $2.96$\,s & $38.1$\,ms & 5\,KB\\
\bottomrule
\end{tabular}
\caption{halo2 formal-ZK backend, Apple~M4. \emph{Cold} regenerates the
deterministic IPA parameters and keys per call; \emph{cached} is the per-proof
cost once they are reused.}
\label{tab:halo2}
\end{table}

\paragraph{Commitment overhead and throughput.} On the default Metal path the
Layer-1 commitment work is $2.6\%$ of inference time at warm-cache decode speed
(${\approx}79$\,tok/s), $0.4\%$ cold. The Layer-2 verifier re-executes roughly
$k/L$ prompt-forward equivalents; $60$ challenges over a $60$-position trace
verify in ${\approx}2$\,s on CPU.

\paragraph{Deterministic backend.} On the 1B model the deterministic backend
decodes at $10.2$\,tok/s (versus ${\approx}80$ on Metal), loads in $2.2$\,s, and
is bit-identical across runs; $20$ real-model challenges verify with exactly
zero deviation. This is the configuration that gives exact token provenance.

\section{Related Work}\label{sec:related}

\emph{zkML.} Systems such as those of Kang et al.~\cite{zkml} and toolchains like
ezkl~\cite{ezkl} prove DNN inference in zero knowledge, typically for smaller
models or restricted layers, and often with a trusted setup. We deliberately do
\emph{not} prove the forward pass; our zero-knowledge component is confined to
the decode (argmax) step, which keeps the strongest guarantee tractable on a
laptop while the cheaper layers cover the rest of the computation.

\emph{Verifiable computation and spot checks.} Re-executing a random subset of a
committed computation echoes probabilistically checkable verification and
Freivalds' randomized checking~\cite{freivalds}; our Layer~2 applies it to
transformer blocks with a Merkle-committed activation trace and a Fiat--Shamir
(or interactive-nonce) challenge.

\emph{Transparency logs.} The append-only, domain-separated Merkle construction
follows Certificate Transparency (RFC~6962)~\cite{rfc6962}; we reuse its tree
shape and second-preimage hardening for the model commitment.

\emph{Proof systems.} We build on FRI-based STARKs~\cite{stark} and the halo2
inner-product argument~\cite{bulletproofs,halo}, with Rescue-Prime~\cite{rescue}
and Poseidon~\cite{poseidon} as circuit-friendly commitments.

\section{Limitations and Future Work}\label{sec:limits}

The strongest guarantee is confined to greedy decoding; committed-seed top-$p$
and temperature proofs are future work. The halo2 backend, though formally
zero-knowledge, is demonstrated only at small vocabularies; a full-vocabulary
formal-ZK decode proof is minutes-to-GB and not yet practical. The STARK's lack
of trace randomization makes it an argument, not formal ZK; a wrapping or
upstream trace randomization would remove the residual confirmation risk. The
commitment binds the quantized bytes and backend, not a hardware attestation:
the deterministic backend gives cross-machine reproducibility for a fixed binary
version, but a malicious runtime substituting the binary is out of scope. Like
all research cryptography, the constructions are unaudited.

\section{Conclusion}\label{sec:conclusion}

\textbf{verifiable-llm} shows that useful, checkable provenance for local LLM
inference is achievable today by declining to solve the hardest version of the
problem and instead composing layered guarantees with honest, quantified
security statements. Its most transferable contribution may be methodological:
it attacks its own weakest layer, measures exactly how far the attack reaches,
and ships a fix rather than a caveat. The result is a system whose credibility
rests on precision and understatement---a stance we believe applied
cryptographic systems should adopt more often.

\begin{thebibliography}{99}
\small
\bibitem{blake3} J.-P. Aumasson, S. Neves, Z. Wilcox-O'Hearn, and J. O'Connor.
\emph{BLAKE3: One Function, Fast Everywhere.} Specification, 2020.
\bibitem{rfc6962} B. Laurie, A. Langley, and E. Kasper. \emph{Certificate
Transparency.} RFC~6962, IETF, 2013.
\bibitem{fiatshamir} A. Fiat and A. Shamir. \emph{How to Prove Yourself:
Practical Solutions to Identification and Signature Problems.} CRYPTO, 1986.
\bibitem{stark} E. Ben-Sasson, I. Bentov, Y. Horesh, and M. Riabzev.
\emph{Scalable, Transparent, and Post-Quantum Secure Computational Integrity.}
IACR ePrint 2018/046.
\bibitem{bulletproofs} B. B\"unz, J. Bootle, D. Boneh, A. Poelstra, P. Wuille,
and G. Maxwell. \emph{Bulletproofs: Short Proofs for Confidential Transactions
and More.} IEEE S\&P, 2018.
\bibitem{halo} S. Bowe, J. Grigg, and D. Hopwood. \emph{Recursive Proof
Composition without a Trusted Setup (Halo).} IACR ePrint 2019/1021.
\bibitem{poseidon} L. Grassi, D. Khovratovich, C. Rechberger, A. Roy, and
M. Schofnegger. \emph{Poseidon: A New Hash Function for Zero-Knowledge Proof
Systems.} USENIX Security, 2021.
\bibitem{rescue} A. Aly, T. Ashur, E. Ben-Sasson, S. Dhooghe, and
A. Szepieniec. \emph{Design of Symmetric-Key Primitives for Advanced
Cryptographic Protocols (Rescue-Prime).} IACR Trans.\ Symmetric Cryptology,
2020.
\bibitem{zkml} D. Kang, T. Hashimoto, I. Stoica, and Y. Sun. \emph{Scaling up
Trustless DNN Inference with Zero-Knowledge Proofs.} arXiv:2210.08674, 2022.
\bibitem{ezkl} ezkl contributors. \emph{ezkl: Zero-Knowledge Inference for Neural
Networks.} \url{https://github.com/zkonduit/ezkl}.
\bibitem{freivalds} R. Freivalds. \emph{Probabilistic Machines Can Use Less
Running Time.} IFIP Congress, 1977.
\bibitem{llama3} Llama Team, AI @ Meta. \emph{The Llama 3 Herd of Models.}
arXiv:2407.21783, 2024.
\end{thebibliography}

\end{document}
